Lie 群告诉我们在同一个时空点上允许怎样旋转内部向量，但局域规范变换允许这个内部基底随位置改变。于是还缺一个新的规则：怎样比较 $x$ 点与 $x+\dd x$ 点的内部向量。规范联络正是这条比较规则；曲率则回答“沿不同路径比较以后，结果是否一致”。在进入联络之前，先用微分形式把反对称指标和边界积分写成一种不依赖坐标排列的语言。

对 $p$ 阶微分形式
\begin{equation*}
 \omega^{(p)}
 =\frac1{p!}\omega_{\mu_1\cdots\mu_p}(x)
 \dd x^{\mu_1}\wedge\cdots\wedge\dd x^{\mu_p},
\end{equation*}
楔积满足 $\dd x^\mu\wedge\dd x^\nu=-\dd x^\nu\wedge\dd x^\mu$。外微分定义为
\begin{equation}
 \dd\omega^{(p)}
 =\frac1{p!}\partial_\nu\omega_{\mu_1\cdots\mu_p}
 \dd x^\nu\wedge\dd x^{\mu_1}\wedge\cdots\wedge\dd x^{\mu_p}.
 \label{eq:exterior-derivative-definition-dp68}
\end{equation}
由偏导交换与楔积反对称性得到两条基本规则：
\begin{equation}
 \boxed{
 \dd^2=0,
 \qquad
 \dd(\alpha^{(p)}\wedge\beta)
 =(\dd\alpha^{(p)})\wedge\beta
 +(-1)^p\alpha^{(p)}\wedge\dd\beta.}
 \label{eq:exterior-calculus-identities-dp68}
\end{equation}
与局域恒等式 $\dd^2=0$ 配套的全局积分关系是 Stokes 定理
\begin{equation}
 \boxed{\displaystyle \int_M\dd\omega=\int_{\partial M}\omega.}
 \label{eq:stokes-forms-dp11}
\end{equation}
它把“体积分变边界项”的计算统一为同一条数学定理。

现在回到局域内部空间。令 $T^a$ 为厄米生成元，矩阵值联络一形式与协变导数定义为
\begin{equation}
 \boxed{
 \mathscr C=\mathscr C_\mu^aT^a\dd x^\mu,
 \qquad
 D=\dd+\ii g_{\rm G}\mathscr C.}
 \label{eq:connection-oneform-dp11}
\end{equation}
$\mathscr C$ 的作用是规定相邻点的内部向量怎样比较。连续作用两次 $D$ 后，路径依赖的最低阶部分定义曲率：
\begin{equation}
 \boxed{
 \mathcal F\equiv\dd\mathscr C+\ii g_{\rm G}\mathscr C\wedge\mathscr C,
 \qquad
 D^2=\ii g_{\rm G}\mathcal F.}
 \label{eq:curvature-form-dp11}
\end{equation}
把二形式展开成分量，
\begin{equation}
 \boxed{
 \mathcal F_{\mu\nu}
 =\partial_\mu\mathscr C_\nu-\partial_\nu\mathscr C_\mu
 +\ii g_{\rm G}[\mathscr C_\mu,\mathscr C_\nu].}
 \label{eq:curvature-components-dp68}
\end{equation}
阿贝尔理论中最后的对易子消失；非阿贝尔场强的额外项直接来自矩阵不交换。

局域基变换 $\Psi'=U\Psi$ 不能改变“平行”的物理含义，因此要求
\begin{equation}
 D'\Psi'=U(D\Psi).
 \label{eq:covariant-derivative-covariance-dp68}
\end{equation}
这唯一固定联络的有限变换律
\begin{equation}
 \boxed{
 \mathscr C'
 =U\mathscr C U^{-1}
 +\frac{\ii}{g_{\rm G}}(\dd U)U^{-1}.}
 \label{eq:finite-gauge-connection-dp11}
\end{equation}
并进一步给出曲率的协变变换
\begin{equation}
 \boxed{\mathcal F'=U\mathcal F U^{-1}.}
 \label{eq:curvature-gauge-covariance-dp68}
\end{equation}
所以 $\Tr(\mathcal F_{\mu\nu}\mathcal F^{\mu\nu})$ 在局域基变换下保持不变。

曲率也可以直接由协变导数的对易子读出：
\begin{equation}
 [D_\mu,D_\nu]=\ii g_{\rm G}\mathcal F_{\mu\nu}.
 \label{eq:covariant-commutator-curvature-dp68}
\end{equation}
把它代入三个协变导数的 Jacobi 恒等式，就得到
\begin{equation}
 \boxed{
 D_\lambda\mathcal F_{\mu\nu}
 +D_\mu\mathcal F_{\nu\lambda}
 +D_\nu\mathcal F_{\lambda\mu}=0.}
 \label{eq:nonabelian-bianchi-dp11}
\end{equation}
这就是非阿贝尔 Bianchi 恒等式。

联络的几何意义最直接地体现在沿曲线输运。对曲线 $C:x^\mu(s)$，平行输运方程为
\begin{equation}
 \frac{\dd\Psi}{\dd s}
 +\ii g_{\rm G}\dot x^\mu(s)\mathscr C_\mu[x(s)]\Psi(s)=0.
 \label{eq:parallel-transport-ode-dp68}
\end{equation}
不同参数点上的矩阵一般不交换，所以有限解必须保留乘法次序：
\begin{equation}
 \boxed{
 U_C(s_f,s_i)
 =\mathcal P\exp\!\left[-\ii g_{\rm G}\int_C\mathscr C\right].}
 \label{eq:wilson-line-general-dp11}
\end{equation}
这个对象称为 Wilson 线。规范变换只旋转它的两个端点，
\begin{equation}
 \boxed{
 U_C' =U(x_f)U_CU^{-1}(x_i).}
 \label{eq:wilson-endpoint-covariance-dp68}
\end{equation}
因此对闭合曲线 $x_f=x_i$ 取迹后，端点旋转在迹中相消，得到规范不变的 Wilson 圈
\begin{equation}
 \boxed{
 W_R(C)=\frac1{d_R}\Tr_R\,
 \mathcal P\exp\!\left[-\ii g_{\rm G}\oint_C\mathscr C\right].}
 \label{eq:wilson-loop-general-dp11}
\end{equation}

曲率为什么可以理解成“局部的不闭合量”，可以直接看一个无穷小矩形。沿四条边依次做平行输运并保留面积阶，
\begin{equation}
 \boxed{
 U_\square
 =I-\ii g_{\rm G}\mathcal F_{\mu\nu}(x)
 \delta x^\mu\delta x^\nu
 +\order(\delta x^3).}
 \label{eq:small-loop-curvature-dp11}
\end{equation}
所以 $\mathcal F_{\mu\nu}$ 正是闭合小回路输运偏离单位矩阵的面积密度。

前面的讨论都是局域的。还需要一个与非 Abel 规范场全局结构直接相关的结果：在四维 Euclidean 时空中，对作用量有限并在无穷远趋于纯规范的 $SU(N)$ 场，可以定义
\begin{equation}
 Q_{\rm top}
 \equiv\frac{g_{\rm G}^2}{8\pi^2}\int\Tr(\mathcal F\wedge\mathcal F).
 \label{eq:pontryagin-form-definition-dp68}
\end{equation}
用 $\Tr(T^aT^b)=\delta^{ab}/2$ 展开分量后，
\begin{equation}
 \boxed{
 Q_{\rm top}
 =\frac{g_{\rm G}^2}{32\pi^2}\int\dd^4 x\,
 \mathcal F_{\mu\nu}^a{}^\star\!\mathcal F^{a\mu\nu}
 \in\mathbb Z.}
 \label{eq:pontryagin-number-dp11}
\end{equation}
这里从微分形式到分量公式只是代数换写；$Q_{\rm top}$ 为整数则还需要无穷远边界条件与整体拓扑定理。也就是说，“归一化系数是多少”可以逐行推导，而“为什么只能取整数”不是局域代数能够单独证明的结论。

Lie 群告诉我们在单个点上允许怎样改变内部基底；联络规定从一个点走到相邻点时怎样比较两个基底；曲率则测量沿两个不同方向依次比较所产生的差异。阿贝尔电磁理论的内部空间是一维相位，矩阵自动交换；非阿贝尔规范理论的内部空间是多维的，比较顺序因此进入场强本身。

